\section{Budget and Resources}
\label{app:resources}

The project was funded by a \href{https://bluedot.org}{Blue Dot Impact} rapid
grant of NZ\$2,542, a Victoria University of Wellington programme contribution
of NZ\$500, and a Victoria University of Wellington AI Centre contribution of
NZ\$1,000, for a total of NZ\$4,042. The Blue Dot Impact rapid grant was the
external award; the other two contributions were internal university support.

\begin{table}[t]
	\centering
	\caption{Project budget: planned versus actual, in New Zealand Dollars}
	\label{tab:budget-detail}
	\begin{tabular}{@{}lrrr@{}}
		\toprule
		Item & Planned & Actual & Difference \\
		\midrule
		\multicolumn{4}{@{}l}{\textbf{Funding received}} \\
		Blue Dot Impact rapid grant & 2,542 & 2,542 & -- \\
		VUW programme contribution & 500 & 500 & -- \\
		VUW AI Centre contribution & 1,000 & 1,000 & -- \\
		\cmidrule(l){2-4}
		\textbf{Total funding} & 4,042 & 4,042 & -- \\
		\midrule
		\multicolumn{4}{@{}l}{\textbf{Expenses}} \\
		GPU compute & 320 & 205 & $-115$ \\
		Participant compensation & 2,625 & 2,925 & $+300$ \\
		Recruitment admin fee & 500 & 0 & $-500$ \\
		Contingency & 689 & 0 & $-689$ \\
		\cmidrule(l){2-4}
		\textbf{Total expenses} & 4,134 & 3,130 & $-1{,}004$ \\
		\midrule
		\textbf{Balance} & $-92$ & 912 & $+1{,}004$ \\
		\bottomrule
	\end{tabular}
\end{table}

Actual expenditure was NZ\$3,130 against a planned NZ\$4,134, leaving a
balance of NZ\$912 (23\% of the funding). The underspend is mostly the
contingency that was not drawn (NZ\$689) and the recruitment admin fee that
was not charged (NZ\$500), partly offset by participant compensation
(NZ\$2,925 against NZ\$2,625). GPU compute was NZ\$205 against NZ\$320.
Participant compensation is an approximate figure of 500 respondents at
NZ\$5.85, with no separate recruitment admin fee. The GPU figure is the
vast.ai bill for 130.9 metered hours, including the 42.7-hour instance left
running by mistake; university GPU time carried no marginal cost.

\subsection{GPU use}
\label{app:gpu-use}

Non-trivial hardware resources were used in this project. The main requirement was GPU compute for the fine-tuning and then serving the models for evaluations and dilemmas. All the GPU compute was done on a single RTX 6000 Pro. Initially a university provided a GPU; however, due to conflicting demands a cloud provider (vast.ai) was used for most of the project. Below is the breakdown of the costs.

\begin{table}[t]
	\centering
	\caption{GPU use by project stage. All hours are wall-clock hours on a
		single RTX PRO 6000. ``Reproduce'' is the cost of rerunning the final
		pipeline from the published code and data; ``Development'' covers
		earlier model sizes, superseded evaluation runs and debugging. Cloud
		hours come from the vast.ai billing record; fine-tuning ran on
		university hardware and was not billed.}
	\label{tab:gpu-use}
	\begin{tabular}{@{}lrrr@{}}
		\toprule
		Stage & Reproduce (h) & Development (h) & Total (h) \\
		\midrule
		Fine-tuning (28 LoRA adapters)   & 3.7  & 16.5 & 20.2 \\
		Evaluation and serving           & 30.4 & 54.6 & 85.0 \\
		Setup and debugging              & --   & 5.0  & 5.0  \\
		Instance left running by mistake & --   & 42.7 & 42.7 \\
		\midrule
		Total                            & 34.1 & 118.8 & 152.9 \\
		\bottomrule
	\end{tabular}

	\smallskip
	\footnotesize{Billed cloud use was 130.9 GPU-hours at \$205.26:
		evaluation 83.3 h (\$137.02), setup 5.0 h (\$9.07) and the accidental
		instance 42.7 h (\$59.17). The other 1.7 h of evaluation and all
		20.2 h of fine-tuning ran on university hardware at no marginal cost.}
\end{table}

\subsection{OpenRouter API use}
\label{app:openrouter-use}

Frontier models were used for the dilemma response judge, the three-judge
response-quality ensemble and the behavioural simulations. The account total
is read directly from the OpenRouter credits endpoint; the split by purpose is
reconstructed from the token counts in the run logs and OpenRouter list prices.

\begin{table}[t]
	\centering
	\caption{OpenRouter API use by purpose. Spend is the account total from
		\texttt{/api/v1/credits}; the split is reconstructed from token counts
		in the run logs and OpenRouter list prices. The unallocated remainder
		is calls whose caches were overwritten or that were made outside the
		repository.}
	\label{tab:openrouter-use}
	\begin{tabular}{@{}lrr@{}}
		\toprule
		Purpose & Spend (USD) & Share \\
		\midrule
		Dilemma response-quality judging         & 56.8 & 57.4\% \\
		Dilemma responses (frontier models and judges) & 5.3 & 5.4\% \\
		Survey quality review                    & 3.0  & 3.0\%  \\
		Frontier model evaluations               & 2.0  & 2.0\%  \\
		Dilemma sourcing and spot tests          & 1.0  & 1.0\%  \\
		Behavioural simulations                  & 0.4  & 0.4\% \\
		Unallocated (superseded passes and retries) & 30.5 & 30.8\% \\
		\midrule
		Total                                    & 99.0 & 100.0\% \\
		\bottomrule
	\end{tabular}

	\smallskip
	\footnotesize{Account total \$99.01 of the \$100 credit purchased, of which
		\$98.82 was spent on the project key. The quality-judging figure
		includes a superseded single-subset pass; the scores retained in the
		final caches cost \$50.6 at list prices. The survey-review, frontier-evaluation
		and sourcing rows are estimated from call counts; the others are
		measured from token logs.}
\end{table}

\section{AI use}

AI agents were used in this project to increase the scope of the project within the limited person-hours. They were used as coding agents to write much of the experimental and analysis code, to search for relevant research, to help proofread the report and to write the LaTeX code for the report. Most of this was done using OpenCode with DeepSeek V4.1 Flash; however, other models were used.

\section{Pre-registration}

A pre-registration was done for both of the Research Questions.

It can be found here: \url{https://aspredicted.org/kd3an7.pdf} \cite{thompsonPreregistrationAligningLLMs2026/10/1515:17}

\subsection{Deviations from the pre-registration}

We added a lower-bound cutoff of four minutes for survey responses, because anything shorter was deemed too low effort. This differs from the pre-registered rule of screening out responses faster than one-third of the median completion time. The screenout was applied live so that the panel provider could replace the respondent.